\section*{第 \ref{ch:b3:gfd-foundations} 章——地球物理流体动力学}

\begin{solution}{exo:b3:gfd-foundations:1}
(a)~$f = 2\Omega\sin\varphi$：\num{7.292e-5}、\num{1.031e-4} 和
\qty{1.263e-4}{s^{-1}}；$2\pi/f$：23.93、16.92 和 \qty{13.82}{h}。
(b)~$U/f$：2.74、1.94 和 \qty{1.58}{km}。(c)~在 \ang{30}，$\sin\varphi =
1/2$，且 $2\pi/f = 2\pi/\Omega$：惯性周期即恒星日，
\qty{23.93}{h}。
\end{solution}

\begin{solution}{exo:b3:gfd-foundations:2}
在 $f = \qty{1.031e-4}{s^{-1}}$ 时：尘卷风 \num{9.7e3}、海风 2.4、
低压 0.15、环流 \num{2.4e-4}。低压与环流接近
地转平衡；尘卷风向哪一侧旋转都可以，对地球自转毫无察觉，
而海风要经过一天才会感受到它。
\end{solution}

\begin{solution}{exo:b3:gfd-foundations:3}
$\partial u_g/\partial z = (g/fT)\,|\partial T/\partial y| = 9.81\times
\num{5e-6}/(\num{1.031e-4}\times250) = \qty{1.90e-3}{s^{-1}}$；在
\qty{10}{km} 范围内，西风增加 \qty{19}{m/s}。
\end{solution}

\begin{solution}{exo:b3:gfd-foundations:4}
$E = \num{50e-8}$：$N = \qty{2.21e-3}{s^{-1}}$，周期为 \qty{47}{min}，
每小时 \num{1.27} 周。$E = \num{1000e-8}$：$N = \qty{9.90e-3}{s^{-1}}$，周期为
\qty{10.6}{min}，每小时 5.7 周。$E = \num{e-8}$：$N =
\qty{3.13e-4}{s^{-1}}$，周期 \qty{5.6}{h}，每小时 0.18 周。
\end{solution}

\begin{solution}{exo:b3:gfd-foundations:5}
(a)~$W \sim UH/L = \qty{0.1}{m/s}$，且 $\mathrm Dw/\mathrm Dt \sim UW/L =
\qty{e-6}{m/s^2}$，为 $g$ 的 \num{1.0e-7}。压强项 $fUL/H =
\qty{0.103}{m/s^2}$，较前者大 \num{1e5} 倍：该比值 \num{9.7e-6} 为
$\mathrm{Ro}(H/L)^2$，其中 $\mathrm{Ro} = 0.097$。(b)~$\tilde fu =
2\Omega\cos\ang{45}\times10 = \qty{1.03e-3}{m/s^2}$，为 $g$ 的 \num{1.05e-4}。
\end{solution}

\begin{solution}{exo:b3:gfd-foundations:6}
$fr = \qty{30}{m/s}$. (a)~$\mathrm{Ro}_g = 0.667$, $V = 40/(1 + \sqrt{3.667})
= \qty{13.7}{m/s}$，比地转速度慢三分之一。(b)~$\mathrm{Ro}_g =
0.167$, $V = 10/(1 + \sqrt{0.333}) = \qty{6.34}{m/s}$. (c)~$V_g \le fr/4 =
\qty{7.5}{m/s}$，则 $V = fr/2 = \qty{15}{m/s}$，为
地转速度的两倍。
\end{solution}

\begin{solution}{exo:b3:gfd-foundations:7}
(a)~$\Omega = 2\pi/6 = \qty{1.047}{rad/s}$, $f = \qty{2.094}{s^{-1}}$.
$\mathrm{Ro} = 0.01/(2.094\times0.05) = 0.095$，且 $\mathrm{Ek} =
\num{e-6}/(2.094\times0.04) = \num{1.19e-5}$：两者均很小，因此流动满足
泰勒–普劳德曼定理，并形成\omterm{thm:b3:gfd-foundations:taylor-proudman}{泰勒柱}。(b)~$0.2/\sqrt{\num{e-6}
\times1.047} = \qty{195}{s}$，\qty{3.3}{min}，转台转了 33 圈。
\end{solution}

\begin{solution}{exo:b3:gfd-foundations:8}
$\partial\bar T/\partial y = -f\,\Delta u/(R\ln(850/250))$，其中
$\ln(850/250) = 1.224$。在 北纬 \ang{30}，$-\num{7.29e-5}\times36.12/(286.9
\times1.224) = \qty{-7.50}{K}$，每 \qty{1000}{km}；在 北纬 \ang{45}，
$-\num{1.031e-4}\times17.18/351.1 = \qty{-5.05}{K}$，每 \qty{1000}{km}。
切变正比于梯度除以 $f$：梯度大 1.5
倍（\ang{30} 处），$f$ 小 1.4 倍，于是切变大 2.1
倍。
\end{solution}

\begin{solution}{exo:b3:gfd-foundations:9}
(a)~$L_R = 0.0113\times11\,000/\num{1.031e-4} = \qty{1206}{km}$;
$0.00357\times1000/\num{9.47e-5} = \qty{37.7}{km}$，小 32 倍。
(b)~$\beta = 2\Omega\cos\varphi/a$：\num{1.62e-11} 和
\qty{1.74e-11}{m^{-1}.s^{-1}}。$\beta L_R/f = 0.19$ 对应大气，0.007
对应海洋。
\end{solution}

\begin{solution}{exo:b3:gfd-foundations:10}
(a)~$\mathrm dw/\mathrm dt = -(r + \mathrm if)w$，因此 $w = U\mathrm
e^{-rt}\mathrm e^{-\mathrm ift}$：速度以速率 $f$ 转向，同时
其大小按 $\mathrm e^{-rt}$ 衰减。位置 $x + \mathrm iy =
U(1 - \mathrm e^{-(r + \mathrm if)t})/(r + \mathrm if)$ 螺旋趋于
点 $U/(r + \mathrm if)$。(b)~$x + \mathrm iy = U(1 - \mathrm e^{-(r +
\mathrm if)t})/(r + \mathrm if) + u_0t$：螺旋被拉长成环。
(c)~$0.2\,\mathrm e^{-1} = \qty{0.074}{m/s}$；中心已漂移 $u_0t =
0.05\times4\times86\,400 = \qty{17.3}{km}$，向东。
\end{solution}

\begin{solution}{exo:b3:gfd-foundations:11}
$f = \qty{9.37e-5}{s^{-1}}$; $\partial u/\partial z = (g/f\rho_0)\,\partial
\rho/\partial y = 9.81\times10^{-5}/(\num{9.37e-5}\times1025) =
\qty{1.02e-3}{s^{-1}}$。取 $u = 0$ 于 \qty{1000}{m} 深处，则表层流为
\qty{1.02}{m/s}，向东：在北半球，洋流
流动时，轻水位于其右侧。
\end{solution}

\begin{solution}{exo:b3:gfd-foundations:12}
(a)~$N^2 = (g/\rho)\,\Delta\sigma_\theta/\Delta z$. 0--\qty{100}{m}:
$9.81\times0.668/(1026.4\times100)$, $N = \qty{7.99e-3}{s^{-1}}$, 4.6
每小时 4.6 周；100--\qty{500}{m}：\qty{3.02e-3}{s^{-1}}，每小时 1.7
周；500--\qty{1000}{m}：\qty{3.24e-3}{s^{-1}}，每小时 1.9 周。
(b)~顶部 \qty{100}{m}：包含位于季节性
密度跃层之上的温暖表层水。100 至 \qty{500}{m} 之间的水体比
下方更均匀。(c)~$E = N^2/g$ 为 \num{650e-8}、\num{93e-8} 和
\qty{107e-8}{m^{-1}}，均在上层海洋的范围内。
\end{solution}

\begin{solution}{pb:b3:gfd-foundations:1}
\textbf{1.} $f = 2\times\num{7.292115e-5}\sin\ang{48} =
\qty{1.084e-4}{s^{-1}}$; $2\pi/f = \qty{16.1}{h}$.

\textbf{2.} 浴缸 \num{3.1e3}；水槽（$f = 4\pi/10 = \qty{1.257}{s^{-1}}$）
0.020；海风 2.3；低压 0.14；环流 \num{2.3e-4}。

\textbf{3.} 浴缸完全感受不到地球自转；海风
只有经过若干小时后才感受得到。水槽、低压与
环流的\omterm{def:b3:gfd-foundations:numbers}{罗斯贝数}都很小，接近地转平衡。

\textbf{4.} 水槽 $\num{e-6}/(1.257\times0.04) = \num{2.0e-5}$；低压 
$10/(\num{1.084e-4}\times10^8) = \num{9.2e-4}$；环流 
$\num{e-3}/(\num{1.084e-4}\times10^6) = \num{9.2e-6}$。

\textbf{5.} $\delta/H = \sqrt{2\nu/(fH^2)} = \sqrt{2\,\mathrm{Ek}}$。水槽 
$\sqrt{2\times\num{e-6}/1.257} = \qty{1.3}{mm}$；空气 $\sqrt{20/\num{1.084e-4}}
= \qty{430}{m}$。

\textbf{6.} 摩擦只作用于若干薄层：在水槽和海洋中，其厚度约为深度的
千分之几，在大气中约为百分之几；这些薄层之上的流动
几乎无摩擦。

\textbf{7.} 空气密度 $\rho = p/(RT) = 101\,290/(286.9\times288.1) =
\qty{1.225}{kg/m^3}$。

\textbf{8.} 地转风 $u_g = (1/\rho f)\,|\partial p/\partial y| = (400/300\,000)/(1.225
\times\num{1.084e-4}) = \qty{10.0}{m/s}$，风从西方吹来，
北侧低压位于左手边。

\textbf{9.} $\mathrm{Ro}_g = 10.04/(\num{1.084e-4}\times600\,000) = 0.154$；
$V = 2\times10.04/(1 + \sqrt{1.618}) = \qty{8.84}{m/s}$，
慢了 \qty{12}{\percent}。

\textbf{10.} $V = 2\times10.04/(1 + \sqrt{0.382}) = \qty{12.4}{m/s}$。最大的 
$V_g$ 为 $fr/4 = \qty{16.3}{m/s}$，对应的梯度 $\rho fV_g =
\qty{2.16e-3}{Pa/m}$，即每 \qty{300}{km} 变化 \qty{6.5}{hPa}。

\textbf{11.} $g/c_p = \qty{9.77}{K/km}$。$N^2 = (9.81/255.7)\times(9.77 -
6.49)\times10^{-3} = \qty{1.26e-4}{s^{-2}}$，$N = \qty{0.0112}{s^{-1}}$，其
周期为 \qty{9.3}{min}。

\textbf{12.} 当 $\Gamma = 0$ 时，$N^2 = (g/T)(g/c_p) = g^2/(c_pT) =
\qty{4.43e-4}{s^{-2}}$，$N = \qty{0.0210}{s^{-1}}$，为
对流层值的 1.9 倍。

\textbf{13.} $L_R = 0.0112\times10^4/\num{1.084e-4} = \qty{1035}{km}$，恰与
低压的尺度相当。

\textbf{14.} $g' = \qty{0.0147}{m/s^2}$；$\sqrt{g'h} = \qty{0.54}{m/s}$；
$L_R = 0.542/\num{1.084e-4} = \qty{5.0}{km}$。

\textbf{15.} 是的：湖泊宽 \qty{12}{km}，超过其
罗斯贝半径的两倍，因此其内部波与水流
受科里奥利力作用而发生偏转。

\textbf{16.} $(249.4 - 253.0)/600 = \qty{-6.0}{K}$，每 \qty{1000}{km}。

\textbf{17.} $\ln(1013/265) = 1.3409$。南方：
$286.9\times253.0\times1.3409/9.81 = \qty{9922}{m}$；北方 \qty{9781}{m}。
\qty{265}{hPa} 等压面在北方低 \qty{141}{m}，
相对于下方等压面的斜率为 \num{2.35e-4}。

\textbf{18.} $\Delta u = (g/f)\times\text{斜率} = 9.81\times\num{2.35e-4}/
\num{1.084e-4} = \qty{21.3}{m/s}$。

\textbf{19.} $(R/f)\ln(p_b/p_t)\,|\partial\bar T/\partial y| =
(286.9/\num{1.084e-4})\times1.3409\times\num{6.0e-6} = \qty{21.3}{m/s}$，结果
相同：两种计算本质上是同一回事。

\textbf{20.} \omterm{def:b3:gfd-foundations:baroclinic}{斜压}：温度（进而密度）沿
等压面变化。若为正压，则风速不随高度变化：
\qty{10.0}{m/s}。

\textbf{21.} $(g/f\bar T)\,|\partial T/\partial y|\,\Delta z = 9.81\times
\num{6.0e-6}\times10^4/(\num{1.084e-4}\times255.7) = \qty{21.2}{m/s}$，
误差在 \qty{0.3}{\percent} 以内。

\textbf{22.} $10.04 + (286.9/\num{1.084e-4})\times\num{6.0e-6}\times
\ln(1013/541) = 10.04 + 9.96 = \qty{20.0}{m/s}$。

\textbf{23.} $\partial u/\partial z = 21.3/10^4 = \qty{2.13e-3}{s^{-1}}$；
$\mathrm{Ri} = \num{1.26e-4}/\num{4.54e-6} = 28$，远大于 0.25：剪切
并不会使自由对流层产生湍流。

\textbf{24.} 奥林的风速 \qty{31}{m/s}，强于
北纬 \ang{45} 处的纬向平均，而弱于北纬 \ang{30} 处的纬向平均。极锋急流南北摆动，
任一时刻其梯度及风速都集中在一条带上，
而沿纬圈取平均则使其分散开来。

\textbf{25.} $10.0 + 21.3$：草原上方 \qty{10}{km} 处的地转风为
\textbf{\qty[text-series-to-math=true]{31}{m/s}}，风从西方吹来
（按计算为 \qty{31.3}{m/s}）。
\end{solution}
